\documentclass[12pt,a4paper]{article}
\usepackage[utf8]{inputenc}
\usepackage[T2A]{fontenc}
\usepackage[russian]{babel}
\usepackage{amsmath,amssymb}
\usepackage{hyperref}
\title{ГРА-Обнуленка: Математический аппарат и методология создания гениального кинематографа}
\author{Исследовательская группа GRA}
\date{\today}
\begin{document}
\maketitle

\begin{abstract}
В работе формализуется понятие \textbf{ГРА-обнуленки} --- рекурсивного процесса аннигиляции целевых иерархий персонажа, приводящего к катарсическому очищению зрителя. Показано, что экспоненциальная сложность обнуленки лежит в основе шедевральности классических фильмов, а современные спецэффектные блокбастеры утрачивают эту глубину. Предложены алгебраические операции над обнуленками, методы их синтеза в сценарии и режиссуре. Приводится практический алгоритм построения фильма как суммы обнуленок с заданными катарсическими свойствами.
\end{abstract}

\section{Введение}
Наблюдение за эволюцией детективного жанра показывает, что фильмы до 2000 года обладают особой ``шедевральностью'', достигаемой через драматургический катарсис, тогда как современное кино всё чаще подменяет его визуальными эффектами. С точки зрения Глубинной Рекурсивной Архитектоники (GRA), причина заключается в наличии или отсутствии \textbf{обнуленки высшей иерархии целей} --- процесса, при котором ложные мотивации персонажа последовательно разрушаются, обнажая экзистенциальный ноль, из которого рождается подлинное очищение.

\section{Формализация обнуленки}
Пусть субъект $S$ описывается иерархическим деревом целей $\mathcal{H} = (V, E, r)$, где $r$ --- корневая (высшая) цель. Обнуленкой $\mathcal{N}$ называется оператор, удовлетворяющий:
\[
\mathcal{N}(\mathcal{H}) = \lim_{n\to\infty} F^n(\mathcal{H}),\quad F(V_i) = V_i \setminus \{\text{подцели, не ведущие к } r\},
\]
при этом финальное состояние $\mathcal{N}(\mathcal{H})$ представляет собой чистый аффект, лишённый исходной структуры.

Сложность вычисления $\mathcal{N}$ экспоненциальна по глубине $d$ и коэффициенту ветвления $b$:
\[
T(\mathcal{N}) = O((2b)^d).
\]

\section{Алгебра обнуленок}
Обнуленки можно суммировать, образуя абелеву группу. Сумма $\mathcal{N}_1 + \mathcal{N}_2$ соответствует последовательному или параллельному обнулению, результат коммутативен и ассоциативен. В линейном пространстве обнуленок определены:
\begin{itemize}
\item Взвешенные комбинации: $\alpha\mathcal{N}_1 + \beta\mathcal{N}_2$ --- регулировка интенсивности катарсиса.
\item Скалярное произведение: $\langle \mathcal{N}_1, \mathcal{N}_2 \rangle$ --- мера ортогональности (независимости) катарсических векторов.
\end{itemize}

Фундаментальное тождество интерсубъективного синтеза:
\[
\neg S \oplus S' \longrightarrow S'',
\]
где $\neg S$ --- результат обнуления субъекта $S$, $\oplus$ --- оператор слияния, а $S''$ --- новый субъект-зритель, прошедший катарсис.

\section{Применение в сценарном деле}
Сценарий гениального фильма строится как последовательность обнуленок всё более высокого порядка. Кульминация есть обнуление корневой цели, что приводит к катарсису. Практический алгоритм:
\begin{enumerate}
\item Определить исходную иерархию целей протагониста.
\item Выявить ложные цели, подлежащие обнулению.
\item Спроектировать сцены-обнуленки в порядке возрастания глубины.
\item Завершить полной аннигиляцией высшей цели, оставив зрителя в состоянии очищения.
\end{enumerate}

\section{Заключение}
Метод ГРА-обнуленок даёт сценаристам и режиссёрам точный инструмент для создания глубокого, преобразующего искусства. Возвращение к культуре обнуления позволит преодолеть кризис современного зрелищного кино.

\end{document}
